\documentclass[10pt,a4paper]{amsart}
\usepackage[margin=19mm]{geometry}
\usepackage{amsmath,amssymb,mathtools}
\usepackage[scr=rsfs,cal=euler]{mathalfa}
\usepackage{xfrac}
\usepackage[unicode,hidelinks,bookmarks=false]{hyperref}
\newcommand{\ie}{i.e.\ }
\newcommand{\cf}{cf.\ }
\newcommand{\resp}{respectively\ }
\newcommand{\Subsection}{\subsection}
\providecommand{\nobreakdash}{\nobreak\hbox{-}}
\setlength{\emergencystretch}{2em}
\multlinegap0pt
\usepackage{amssymb}
\usepackage[shortlabels]{enumitem}
\usepackage{xcolor}
\newtheorem{lemma}{Lemma}
\newtheorem{theorem}{Theorem}
\newcommand{\E}{\mathbb E}
\newcommand{\Kk}{\mathbf K}
\newcommand{\N}{\mathbb N}
\newcommand{\Pbb}{\mathbb P}
\newcommand{\R}{\mathbb R}
\newcommand{\ind}[1]{\mathbf 1_{\{#1\}}}
\title{Convergence of a Critical Multitype Branching Particle System with Mixed Finite and Infinite Mean Lifetimes}
\author{J.H. Ram\'{i}rez-Gonz\'{a}lez, Fabio Prates Machado}
\date{Source: arXiv:2606.11511v2 (CC BY 4.0)}
\begin{document}
\maketitle

\noindent\emph{Source and licence.} Convergence of a Critical Multitype Branching Particle System with Mixed Finite and Infinite Mean Lifetimes. Version arXiv:2606.11511v2 (CC BY 4.0), distributed by arXiv under the Creative Commons Attribution 4.0 International licence, http://creativecommons.org/licenses/by/4.0/, as recorded in the arXiv licence metadata for this exact version and shown on the abstract page https://arxiv.org/abs/2606.11511v2. Full licence notice: \texttt{licenses/arxiv-2606.11511-CC-BY-4.0.txt}.\par\smallskip

\noindent\textit{Editorial scope.} Counted unit: the article's theorem thm:mixed-renewal (source0.tex lines 1067--1099). The article's complete original proof of it is delivered with it (source0.tex lines 1101--1754, one \texttt{proof} environment of 654 lines). No mathematics is written, repaired or re-derived: the statement and the proof are the article's own lines, in the article's own notation, and no retained line is substituted or re-flowed.  Retained uncounted as the source-local context the proof needs: lines 322--326; lines 814--842; lines 853--856.  Nothing was omitted from any retained span, so the package carries no omission record.  No retained line contains a citation, so the wrapper carries no bibliography and no bibliography entry.  No retained line contains a figure environment, an image-inclusion command or drawing code, so no image is delivered and the package declares no asset dependency.  Editorial formatting only, in addition: an AMS \texttt{amsart} preamble that restates the article's own declarations and macros used by the retained lines, namely the environments \texttt{lemma}, \texttt{theorem} on separate counters, together with the standard packages those lines need.  The retained environments print in this document as Lemma 1, Theorem 1 in order of appearance, whereas the article prints the same items with the numbering of its own full document.  The complete native source of the article as distributed in the arXiv e-print for this exact version is bundled unchanged as \texttt{sources/202556/source0.tex}.
\begin{equation}\label{eq:F1-heavy-tail}
  \overline F_1(t)\sim c_1t^{-\gamma},
  \qquad t\to\infty,
  \qquad c_1>0,\quad 0<\gamma<1.
\end{equation}

\begin{lemma}\label{lem:infinite-mean-renewal-negligibility}
Let \(X\) be a strictly positive, non-arithmetic random variable with distribution function \(F\), renewal measure
\[
  U:=\sum_{n=0}^{\infty}F^{*n},
  \qquad F^{*0}:=\delta_0,
\]
and tail
\[
  \overline F(t):=1-F(t)=\Pbb\{X>t\}
  \sim ct^{-\gamma},
  \qquad t\to\infty,
  \qquad c>0,\quad 0<\gamma<1.
\]
Let \(h:[0,\infty)\to\R\) be locally bounded and Borel measurable, and extend it by zero to \((-\infty,0)\). If
\begin{equation}\label{eq:renewal-small-relative-tail}
  |h(t)|=o\bigl(\overline F(t)\bigr),
  \qquad t\to\infty,
\end{equation}
then
\begin{equation}\label{eq:renewal-negligibility-X}
  \int_{[0,t]}h(t-s)\,U(ds)\longrightarrow0,
  \qquad t\to\infty.
\end{equation}
In particular, for every \(\kappa>\gamma\),
\begin{equation}\label{eq:renewal-power-negligibility-X}
  \int_{[0,t]}(1+t-s)^{-\kappa}\,U(ds)\longrightarrow0,
  \qquad t\to\infty.
\end{equation}
\end{lemma}

\begin{equation}\label{eq:renewal-tail-identity-general}
  \int_{[0,t]}\overline F(t-s)\,U(ds)=1,
  \qquad t\ge0.
\end{equation}

\begin{theorem}\label{thm:mixed-renewal}
Assume that \(M\) is irreducible and stochastic, with \(m_{1,1}>0\), that \(F_i(0)=0\) for every \(i\), that \(F_1\) is non-arithmetic and satisfies \eqref{eq:F1-heavy-tail}, and that
\[
  \int_0^\infty\overline F_i(t)\,dt<\infty,
  \qquad i=2,\dots,K.
\]
Let \(z_1,\dots,z_K\) be real-valued Borel measurable functions, locally bounded on \([0,\infty)\), and extended by zero to \((-\infty,0)\). Assume that there exists \(a\in\R\) such that
\begin{equation}\label{eq:z1-decomp-renewal}
  z_1(t)=a\overline F_1(t)+r_1(t),
  \qquad t\ge0.
\end{equation}
Assume that each function among \(r_1,z_2,\dots,z_K\) satisfies
\begin{equation*}\tag{H}
  |h(t)|=o\bigl(\overline F_1(t)\bigr),
  \qquad t\to\infty.
\end{equation*}
where \(h\) denotes the function under consideration.
Then the renewal system
\begin{equation}\label{eq:mixed-renewal-system-final}
  H_i(t)
  =
  z_i(t)
  +
  \sum_{j=1}^K
  \int_{[0,t]}H_j(t-s)\,dF_{i,j}(s),
  \qquad i=1,\dots,K,
\end{equation}
has a unique locally bounded Borel measurable solution \(H=(H_1,\dots,H_K)\). Moreover,
\begin{equation}\label{eq:mixed-renewal-limit-final}
  \lim_{t\to\infty}H_i(t)=a,
  \qquad i=1,\dots,K.
\end{equation}
\end{theorem}

\begin{proof}
We first establish non-explosion and the local finiteness estimate needed below. Fix \(\theta>0\), and set
\[
  \phi_j(\theta):=
  \int_{[0,\infty)}e^{-\theta s}\,dF_j(s),
  \qquad
  q_\theta:=\max_{1\le j\le K}\phi_j(\theta).
\]
Since \(F_j(0)=0\) for every \(j\) and the type space is finite,
\[
  q_\theta<1.
\]
The Markov--renewal construction and the independence of the lifetime variables give
\[
  \E_i\left(e^{-\theta S_n}\right)
  =
  \E_i\left[
    \prod_{r=0}^{n-1}\phi_{Y_r}(\theta)
  \right]
  \le q_\theta^n.
\]
Since \((S_n)_{n\ge0}\) is nondecreasing, the extended limit
\[
  S_\infty:=\lim_{n\to\infty}S_n
  \in[0,\infty]
\]
exists almost surely. With the convention \(e^{-\theta\infty}=0\),
\[
  e^{-\theta S_n}\longrightarrow e^{-\theta S_\infty}
  \qquad\text{almost surely}.
\]
Since \(0\le e^{-\theta S_n}\le1\), the dominated convergence theorem gives
\[
  \E_i\left(e^{-\theta S_\infty}\right)
  =
  \lim_{n\to\infty}\E_i\left(e^{-\theta S_n}\right)
  =0.
\]
Therefore
\[
  e^{-\theta S_\infty}=0
  \qquad \Pbb_i\text{-almost surely},
\]
which is equivalent to
\[
  S_n\longrightarrow\infty
  \qquad \Pbb_i\text{-almost surely}.
\]
For \(i=1,\dots,K\), define
\[
  \widehat H_i(t):=
  \E_i\left[
    \sum_{n\ge0}
    z_{Y_n}(t-S_n)\ind{S_n\le t}
  \right],
  \qquad t\ge0.
\]
With \(\theta>0\) and \(q_\theta<1\) fixed as above, let \(T>0\) and set
\[
  C_T:=
  \max_{1\le j\le K}\sup_{0\le u\le T}|z_j(u)|.
\]
Since
\[
  \ind{S_n\le t}
  \le e^{\theta(t-S_n)}
  \le e^{\theta T}e^{-\theta S_n},
  \qquad 0\le t\le T,
\]
we have
\[
\begin{split}
  \E_i\left[
    \sum_{n\ge0}
    |z_{Y_n}(t-S_n)|\ind{S_n\le t}
  \right]
  &\le
  C_Te^{\theta T}
  \sum_{n\ge0}\E_i\left(e^{-\theta S_n}\right)\\
  &\le
  C_Te^{\theta T}\sum_{n\ge0}q_\theta^n
  =
  \frac{C_Te^{\theta T}}{1-q_\theta}
  <\infty.
\end{split}
\]
Thus the series defining \(\widehat H_i(t)\) is absolutely integrable. Since the corresponding expected partial sums are Borel measurable, \(\widehat H_i\) is Borel measurable, and the preceding bound shows that it is locally bounded. Conditioning on the first holding time and the first transition gives
\[
\begin{split}
  \widehat H_i(t)
  &=
  z_i(t)
  +
  \sum_{j=1}^K m_{i,j}\int_{[0,t]}
  \E_j\left[
    \sum_{n\ge0}
    z_{Y_n}(t-s-S_n)\ind{S_n\le t-s}
  \right]dF_i(s)\\
  &=
  z_i(t)
  +
  \sum_{j=1}^K
  \int_{[0,t]}\widehat H_j(t-s)\,dF_{i,j}(s).
\end{split}
\]
Hence \(\widehat H=(\widehat H_1,\dots,\widehat H_K)\) is a locally bounded Borel measurable solution of \eqref{eq:mixed-renewal-system-final}.

To prove uniqueness, let \(H\) and \(\widetilde H\) be two locally bounded Borel measurable solutions of \eqref{eq:mixed-renewal-system-final}, and set
\[
  D_i(t):=H_i(t)-\widetilde H_i(t),
  \qquad i=1,\dots,K.
\]
Then \(D=(D_1,\dots,D_K)\) satisfies the homogeneous system
\[
  D_i(t)
  =
  \sum_{j=1}^K
  \int_{[0,t]}D_j(t-s)\,dF_{i,j}(s).
\]
Since \(S_1=\tau_0\) and
\[
  \Pbb_i\{Y_1=j,\ S_1\in ds\}
  =
  m_{i,j}\,dF_i(s)
  =
  dF_{i,j}(s),
\]
the preceding equation is equivalently
\[
  D_i(t)
  =
  \E_i\left[
    D_{Y_1}(t-S_1)\ind{S_1\le t}
  \right].
\]
Repeated application of this identity, using the Markov--renewal property after each transition, gives, for every \(n\ge1\),
\[
  D_i(t)
  =
  \E_i\left[
    D_{Y_n}(t-S_n)\ind{S_n\le t}
  \right].
\]
Fix \(T>0\) and set
\[
  M_T:=
  \max_{1\le j\le K}\sup_{0\le u\le T}|D_j(u)|<\infty.
\]
For \(0\le t\le T\),
\[
  |D_i(t)|
  \le
  M_T\Pbb_i\{S_n\le t\}
  \le
  M_T\Pbb_i\{S_n\le T\}.
\]
Since \(S_n\to\infty\) almost surely,
\[
  \Pbb_i\{S_n\le T\}\longrightarrow0,
  \qquad n\to\infty.
\]
Letting \(n\to\infty\) in the preceding bound yields \(D_i(t)=0\). Since \(T>0\) is arbitrary, \(D_i(t)=0\) for every \(t\ge0\) and every \(i\in\Kk\). Consequently, \(\widehat H\) is the unique locally bounded Borel measurable solution of \eqref{eq:mixed-renewal-system-final}. Denoting this solution by \(H\), we obtain
\begin{equation}\label{eq:renewal-representation-final}
  H_i(t)
  =
  \E_i\left[
    \sum_{n\ge0}
    z_{Y_n}(t-S_n)\ind{S_n\le t}
  \right],
  \qquad t\ge0,\quad i=1,\dots,K.
\end{equation}

We next analyse the return cycles to type \(1\). Since \(M\) is irreducible on the finite state space \(\Kk\) and \(m_{1,1}>0\), there exist \(r\in\N\) and \(p_0>0\) such that
\[
  \inf_{i\in\Kk}\Pbb_i\{\sigma_1\le r\}\ge p_0.
\]
By the Markov property,

\begin{equation}\label{eq:return-exp-tail-final}
  \Pbb_i\{\sigma_1>nr\}\le(1-p_0)^n,\qquad n\ge0,\quad i\in\Kk.
\end{equation}
Consequently,
\[
  \Pbb_i\{\sigma_1=\infty\}=\lim_{n\to\infty}\Pbb_i\{\sigma_1>nr\}=0,\qquad \E_i(\sigma_1)=\sum_{m\ge0}\Pbb_i\{\sigma_1>m\}\le r\sum_{n\ge0}\Pbb_i\{\sigma_1>nr\}\le\frac{r}{p_0}<\infty.
\]

Let
\[
  \overline F_C(t):=\Pbb_1\{C>t\}.
\]
We claim that
\begin{equation}\label{eq:cycle-tail-final}
  \overline F_C(t)
  \sim\overline F_1(t)
  \sim c_1t^{-\gamma},
  \qquad t\to\infty.
\end{equation}
Write
\[
  C=\tau_0+E,
  \qquad
  E:=\sum_{n=1}^{\sigma_1-1}\tau_n,
\]
with the convention that \(E=0\) when \(\sigma_1=1\). Before the return to type \(1\), the chain visits only types in \(\{2,\dots,K\}\). Set
\[
  \mu_i:=\int_0^\infty\overline F_i(s)\,ds,\quad i=2,\dots,K,
  \qquad
  \mu_*:=\max_{2\le i\le K}\mu_i<\infty.
\]
Since \(M\) is stochastic, it is the transition matrix of the embedded type chain \((Y_n)_{n\ge0}\). By construction, conditionally on this chain, \(\tau_n\) has distribution \(F_{Y_n}\). Moreover, \(Y_n\ne1\) for \(1\le n<\sigma_1\). Therefore, by conditional expectation and Tonelli's theorem,
\[
  \E_1(E)
  =
  \E_1\left[\sum_{n=1}^{\sigma_1-1}\tau_{n}\right]
  \le
  \mu_*\E_1(\sigma_1-1)
  <\infty.
\]
On the event \(\{E>t\}\), one has \(t\le E\), and therefore
\[
  t\Pbb_1\{E>t\}
  =
  \E_1\bigl(t\ind{E>t}\bigr)
  \le
  \E_1\bigl(E\ind{E>t}\bigr).
\]
Moreover, \(\ind{E>t}\to0\) almost surely and
\[
  0\le E\ind{E>t}\le E.
\]
Since \(\E_1(E)<\infty\), dominated convergence yields
\[
  \E_1\bigl(E\ind{E>t}\bigr)\longrightarrow0.
\]
Consequently,
\[
  t\Pbb_1\{E>t\}\longrightarrow0,
\]
or equivalently, \(\Pbb_1\{E>t\}=o(t^{-1})\). Since \(0<\gamma<1\) and \(\overline F_1(t)\sim c_1t^{-\gamma}\), we have \(t^{-1}=o(\overline F_1(t))\), and therefore
\[
  \Pbb_1\{E>t\}
  =
  o\bigl(\overline F_1(t)\bigr).
\]


Since \(E\ge0\),
\[
  \overline F_C(t)\ge\overline F_1(t).
\]
For \(\delta\in(0,1)\),
\[
  \{C>t\}
  \subseteq
  \{\tau_0>(1-\delta)t\}\cup\{E>\delta t\},
\]
and therefore
\[
  \overline F_C(t)
  \le
  \overline F_1((1-\delta)t)+\Pbb_1\{E>\delta t\}.
\]
Dividing by \(\overline F_1(t)\) gives
\[
  \limsup_{t\to\infty}
  \frac{\overline F_C(t)}{\overline F_1(t)}
  \le(1-\delta)^{-\gamma}.
\]
Letting \(\delta\downarrow0\) proves \eqref{eq:cycle-tail-final}. In particular,
\[
  \E_1(C)=\infty.
\]

The distribution of \(C\) is non-arithmetic. Indeed, suppose that
\[
  \Pbb_1\{C\in h\N_0\}=1
\]
for some \(h>0\), and let \(\mathcal A:=\{Y_1=1\}\). Since \(m_{1,1}>0\),
\[
  \Pbb_1(\mathcal A)=m_{1,1}>0.
\]
On \(\mathcal A\), one has \(\sigma_1=1\) and \(C=\tau_0\). Since \(\tau_0\) is independent of \(\mathcal A\),
\[
\begin{split}
  0
  &=
  \Pbb_1\{C\notin h\N_0\}\ge
  \Pbb_1\bigl(\mathcal A,\tau_0\notin h\N_0\bigr)=
  \Pbb_1(\mathcal A)\Pbb_1\{\tau_0\notin h\N_0\}.
\end{split}
\]
It follows that \(F_1\) is arithmetic, contradicting the hypothesis.

Let \(h\) be a locally bounded Borel function on \([0,\infty)\), extended by zero to \((-\infty,0)\), and suppose that \(h\) satisfies \textup{(H)}. We claim that
\begin{equation}\label{eq:basic-negligible-final}
  (U_C*h)(t):=
  \int_{[0,t]}h(t-s)\,U_C(ds)
  \longrightarrow0,
  \qquad t\to\infty.
\end{equation}
Indeed, by condition \textup{(H)} and \eqref{eq:cycle-tail-final},
\[
  \frac{|h(t)|}{\overline F_C(t)}
  =
  \frac{|h(t)|}{\overline F_1(t)}
  \frac{\overline F_1(t)}{\overline F_C(t)}
  \longrightarrow0.
\]
Hence Lemma~\ref{lem:infinite-mean-renewal-negligibility}, applied to the return-cycle distribution \(F_C\), proves \eqref{eq:basic-negligible-final}. Moreover, \(U_C*h\) is bounded on \([0,\infty)\): it is bounded on compact intervals because \(U_C\) is finite on compact intervals and \(h\) is locally bounded, and it is bounded for large arguments because it converges to zero.

For the process started from type \(1\), define
\begin{equation}\label{eq:cycle-reward-final}
  q(t):=
  \E_1\left[
    \sum_{n=0}^{\sigma_1-1}
    z_{Y_n}(t-S_n)\ind{S_n\le t}
  \right].
\end{equation}



For \(n\ge0\), let
\[
  \mathcal G_n
  :=
  \sigma\bigl(
    Y_0,\ldots,Y_n,
    \tau_0,\ldots,\tau_{n-1}
  \bigr),
\]
with the list of holding times understood to be empty when \(n=0\), and
let
\[
  \mathcal G_\infty
  :=
  \sigma\left(\bigcup_{n\ge0}\mathcal G_n\right).
\]
Then each \(\sigma_k\) is a stopping time with respect to
\((\mathcal G_n)_{n\ge0}\). We write
\[
  \mathcal G_{\sigma_k}
  :=
  \left\{
    A\in\mathcal G_\infty:
    A\cap\{\sigma_k\le n\}\in\mathcal G_n
    \text{ for every }n\ge0
  \right\}
\]
for the corresponding stopping-time sigma-field.

Let
\[
  R_k:=S_{\sigma_k},
  \qquad k\ge0.
\]
Since \(Y_{\sigma_k}=1\), the strong Markov property at \(\sigma_k\)
implies that the increments \(R_{k+1}-R_k\) are independent and
identically distributed with law \(F_C\). Hence
\[
  U_C(\cdot)
  =
  \sum_{k\ge0}\Pbb_1\{R_k\in\cdot\}.
\]

For \(k\ge0\), define the reward accumulated during the \(k\)-th return
cycle by
\[
  Q_k(t)
  :=
  \sum_{n=\sigma_k}^{\sigma_{k+1}-1}
  z_{Y_n}(t-S_n)\ind{S_n\le t}.
\]
The cycle intervals partition the nonnegative integers, and the local
absolute-integrability estimate established above gives
\[
\begin{split}
  \sum_{k\ge0}\E_1|Q_k(t)|
  &\le
  \E_1\left[
    \sum_{n\ge0}
    |z_{Y_n}(t-S_n)|\ind{S_n\le t}
  \right]
  <\infty.
\end{split}
\]
Consequently,
\[
  H_1(t)=\sum_{k\ge0}\E_1[Q_k(t)].
\]

Set
\[
  \rho_k:=\sigma_{k+1}-\sigma_k,\qquad
  \widetilde Y_m^{(k)}:=Y_{\sigma_k+m},\qquad
  \widetilde S_m^{(k)}:=S_{\sigma_k+m}-R_k.
\]
Then
\[
\begin{split}
  Q_k(t)
  &=
  \ind{R_k\le t}
  \sum_{m=0}^{\rho_k-1}
  z_{\widetilde Y_m^{(k)}}
  \bigl((t-R_k)-\widetilde S_m^{(k)}\bigr)
  \ind{\widetilde S_m^{(k)}\le t-R_k}.
\end{split}
\]
Conditional on \(\mathcal G_{\sigma_k}\), the shifted sequence
\[
  \left(
    \widetilde Y_m^{(k)},
    \widetilde S_m^{(k)}
  \right)_{m\ge0}
\]
has the same law as the original Markov--renewal sequence under
\(\Pbb_1\). Since \(R_k\) is \(\mathcal G_{\sigma_k}\)-measurable, the
definition of \(q\) yields
\[
  \E_1\left[
    Q_k(t)\mid\mathcal G_{\sigma_k}
  \right]
  =
  q(t-R_k)\ind{R_k\le t}.
\]
Therefore,
\begin{equation}\label{eq:H1-cycle-final}
H_1(t)=
  \sum_{k\ge0}
  \E_1\left[
    q(t-R_k)\ind{R_k\le t}
  \right]=
  \sum_{k\ge0}
  \int_{[0,t]}q(t-s)\,\Pbb_1\{R_k\in ds\}=
  \int_{[0,t]}q(t-s)\,U_C(ds).
\end{equation}

The term corresponding to \(n=0\) in \eqref{eq:cycle-reward-final} is
\[
  z_1(t)=a\overline F_1(t)+r_1(t).
\]
For the remaining terms, the definition of \(\sigma_1\) as the first
return time to type \(1\) implies that
\[
  Y_n\in\{2,\dots,K\},
  \qquad 1\le n<\sigma_1.
\]
Accordingly, for \(j=2,\dots,K\), define the finite measure
\[
  \ell_j(B):=
  \E_1\left[
    \sum_{n=1}^{\sigma_1-1}
    \ind{Y_n=j,\ S_n\in B}
  \right]
\]
for every Borel set \(B\subset[0,\infty)\). Each \(\ell_j\) is finite,
since
\[
  \ell_j([0,\infty))
  \le
  \E_1(\sigma_1)
  <\infty.
\]

For every nonnegative Borel function \(g\), the definition of the mean
counting measure \(\ell_j\) and Tonelli's theorem give
\[
  \int_{[0,\infty)}g(s)\,\ell_j(ds)
  =
  \E_1\left[
    \sum_{n=1}^{\sigma_1-1}
    g(S_n)\ind{Y_n=j}
  \right].
\]
By applying this identity to the positive and negative parts, it also
holds for every Borel function \(g\) satisfying
\[
  \int_{[0,\infty)}|g(s)|\,\ell_j(ds)<\infty.
\]

For fixed \(t\ge0\), set
\[
  g_{j,t}(s)
  :=
  z_j(t-s)\ind{0\le s\le t}.
\]
Since \(z_j\) is locally bounded and \(\ell_j\) is finite,
\[
\begin{split}
  \int_{[0,\infty)}|g_{j,t}(s)|\,\ell_j(ds)
  &\le
  \sup_{0\le u\le t}|z_j(u)|\,
  \ell_j([0,t])
  <\infty.
\end{split}
\]
Consequently,
\[
\begin{split}
  &\E_1\left[
    \sum_{n=1}^{\sigma_1-1}
    z_j(t-S_n)\ind{Y_n=j,\ S_n\le t}
  \right]\\
  &\qquad=
  \int_{[0,t]}z_j(t-s)\,\ell_j(ds).
\end{split}
\]
Summing over \(j=2,\dots,K\), we obtain
\begin{equation}\label{eq:q-decomposition-final}
  q(t)
  =
  a\overline F_1(t)+r_1(t)
  +
  \sum_{j=2}^K
  \int_{[0,t]}z_j(t-s)\,\ell_j(ds).
\end{equation}
Substituting \eqref{eq:q-decomposition-final} into \eqref{eq:H1-cycle-final} and applying Fubini's theorem gives
\begin{equation}\label{eq:H1-conv-decomp-final}
\begin{split}
  H_1(t)
  &=
  a\int_{[0,t]}\overline F_1(t-s)\,U_C(ds)
  +(U_C*r_1)(t)+
  \sum_{j=2}^K
  \int_{[0,t]}(U_C*z_j)(t-s)\,\ell_j(ds).
\end{split}
\end{equation}
For fixed \(t\), the required absolute integrability follows from
\[
  \sup_{0\le u\le t}|z_j(u)|\,
  U_C([0,t])\,\ell_j([0,t])
  <\infty.
\]
By \eqref{eq:basic-negligible-final},
\[
  (U_C*r_1)(t)\longrightarrow0,
  \qquad
  (U_C*z_j)(t)\longrightarrow0.
\]
Each \(U_C*z_j\) is globally bounded. Hence, after extending it by zero to \((-\infty,0)\), dominated convergence with respect to the finite measure \(\ell_j\) yields
\[
  \int_{[0,t]}(U_C*z_j)(t-s)\,\ell_j(ds)
  \longrightarrow0.
\]
Therefore
\begin{equation}\label{eq:H1-main-reduced-final}
  H_1(t)
  =
  a\int_{[0,t]}\overline F_1(t-s)\,U_C(ds)+o(1).
\end{equation}

Set
\[
  d(t):=\overline F_1(t)-\overline F_C(t).
\]
By \eqref{eq:cycle-tail-final},
\[
  |d(t)|=o\bigl(\overline F_C(t)\bigr).
\]
Lemma~\ref{lem:infinite-mean-renewal-negligibility}, applied to \(C\) and \(d\), gives
\[
  (U_C*d)(t)\longrightarrow0.
\]
On the other hand, the elementary renewal identity \eqref{eq:renewal-tail-identity-general}, applied to \(C\), gives
\[
  \int_{[0,t]}\overline F_C(t-s)\,U_C(ds)=1.
\]
Consequently,
\begin{equation}\label{eq:F1-tail-renewal-final}
  \int_{[0,t]}\overline F_1(t-s)\,U_C(ds)
  =
  1+(U_C*d)(t)
  \longrightarrow1.
\end{equation}
Combining \eqref{eq:H1-main-reduced-final} and \eqref{eq:F1-tail-renewal-final}, we obtain
\[
  H_1(t)\longrightarrow a.
\]


Finally, fix \(i\in\{2,\dots,K\}\), and under \(\Pbb_i\) set
\[
  T_i:=S_{\sigma_1}.
\]
By \eqref{eq:return-exp-tail-final},
\[
  \Pbb_i\{\sigma_1<\infty\}=1,
  \qquad
  \E_i(\sigma_1)<\infty,
\]
and hence \(T_i<\infty\) almost surely. Define
\[
  b_i(t):=
  \E_i\left[
    \sum_{n=0}^{\sigma_1-1}
    z_{Y_n}(t-S_n)\ind{S_n\le t}
  \right].
\]
For \(j=2,\dots,K\), define
\[
  \nu_{i,j}(B):=
  \E_i\left[
    \sum_{n=0}^{\sigma_1-1}
    \ind{Y_n=j,\ S_n\in B}
  \right]
\]
for every Borel set \(B\subset[0,\infty)\). Then
\(
  \nu_{i,j}([0,\infty))
  \le
  \E_i(\sigma_1)
  <\infty
\)
and
\[
  b_i(t)
  =
  \sum_{j=2}^K
  \int_{[0,t]}z_j(t-s)\,\nu_{i,j}(ds).
\]
By condition \textup{(H)} and \(\overline F_1(t)\to0\), each \(z_j\), \(j\ge2\), tends to zero. Since each \(z_j\) is locally bounded, it is therefore bounded on \([0,\infty)\). Dominated convergence gives
\begin{equation}\label{eq:prehit-zero-final}
  b_i(t)\longrightarrow0.
\end{equation}

Extend \(H_1\) by zero to \((-\infty,0)\). Splitting \eqref{eq:renewal-representation-final} at the Markov-chain stopping time \(\sigma_1\) and using the strong Markov property of the Markov--renewal sequence gives
\begin{equation}\label{eq:hit-decomp-final}
  H_i(t)
  =
  b_i(t)
  +
  \E_i\left[
    H_1(t-T_i)\ind{T_i\le t}
  \right].
\end{equation}
Since \(H_1\) is locally bounded and converges to \(a\), it is globally bounded. Since \(T_i<\infty\) almost surely,
\[
  H_1(t-T_i)\ind{T_i\le t}
  \longrightarrow a
  \qquad\text{almost surely}.
\]
Dominated convergence therefore gives
\[
  \E_i\left[
    H_1(t-T_i)\ind{T_i\le t}
  \right]
  \longrightarrow a.
\]
Together with \eqref{eq:prehit-zero-final} and \eqref{eq:hit-decomp-final}, this proves
\[
  H_i(t)\longrightarrow a,
  \qquad i=1,\dots,K.
\]
\end{proof}

\end{document}
